
\twocolumn[
\begin{@twocolumnfalse}
\maketitle
\vspace{-1.4em}
\begin{abstract}
Dense retrieval provides an effective similarity primitive, while structured search efficiently handles exact filters. Many useful constraints---such as \emph{minimalist}, \emph{office-appropriate}, or \emph{not sporty}---are latent semantic predicates: they are semantic judgments rather than stored catalog facts, and they are awkward to express as one query--item similarity. We study a \emph{compile once, execute cheaply} design in which expensive semantic supervision is moved offline, items share one low-bit representation, and each latent predicate is compiled into a reusable search-time program.

The original Random Semantic Algebra (\rsa) instantiation uses random orthogonal projections, low-bit quantization, and sparse boosted lookup tables. We extend the study with a centered binary compiler: a 384-dimensional item is stored as 384 sign bits plus two per-item reconstruction scalars (56 bytes/item), while a supervised linear semantic head is quantized to four-bit weights and executed through four packed bit planes (about 216 bytes/predicate). We compare this design against sparse 1/2/4-bit \rsa programs, zero-shot MiniLM directions, a full-precision linear semantic proxy, and a strong PQ64 baseline that compiles the same linear proxy into 64 product-quantization lookup tables.

On 44,072 fashion products, with MiniLM titles as the search-side representation and CLIP images as an independent teacher, we evaluate three strict fit/calibration/test splits and 200 compound queries per split. At 20\% candidate retention, dense retrieval reaches 0.207 recall and 0.425 purity; 4-bit sparse \rsa reaches 0.330/0.571; the centered binary int4 compiler reaches 0.337/0.578; PQ64 reaches 0.352/0.594; and the FP32 semantic proxy reaches 0.363/0.606. The binary compiler therefore recovers 83\% of the incremental FP32 recall over dense retrieval while using 56 bytes/item and about 216 bytes/predicate; PQ64 recovers 93\% but requires about 65.5\,KB/predicate. Quantizing the binary predicate from FP32 to int4 changes recall by only 0.00047. A single-threaded Rust benchmark over the actual held-out codes and learned programs reaches 1.98 million candidates/s for eight simultaneous binary predicates at a 100k candidate pool, versus 2.04M/s for PQ64 and 0.41M/s for FP32. The strongest supported result is thus not that random 4-bit codes are uniquely optimal, but that semantic supervision can be compiled into very small, composable programs over a shared low-bit substrate with a favorable quality--memory--throughput trade-off.
\end{abstract}
\vspace{0.7em}
\end{@twocolumnfalse}
]

\section{Introduction}
Modern search engines expose two mature computational languages. Dense retrieval expresses semantic similarity through inner products or cosine similarity. Structured retrieval expresses exact predicates---brand, size, category, availability---through inverted indexes, bitmaps, or relational filters. A large class of useful constraints is neither: \emph{minimalist}, \emph{office-appropriate}, \emph{technical}, \emph{quiet luxury}, and explicit negations such as \emph{not sporty} are semantic judgments rather than stored facts.

A dense query vector can absorb these words into one score, but logical composition becomes implicit. At the other extreme, a large VLM or LLM can evaluate each predicate directly, but per-candidate model inference inside a retrieval loop is expensive. Recent semantic-operator systems and semantic-filter cascades make this execution problem explicit \citep{patel2025lotus,kim2026semanticfilter,dong2026compilation}.

We study a third design point: \emph{compile once, execute cheaply}. Items are mapped once into a concept-independent low-bit code. Each semantic concept is learned offline as a small executable program over that code. Query understanding then selects and composes already-compiled predicates, leaving ANN retrieval, exact filters, and downstream ranking unchanged.

\begin{figure*}[t]
\centering
\begin{tikzpicture}[node distance=8mm and 9mm, >=Latex, every node/.style={font=\small}]
\tikzstyle{box}=[draw, rounded corners=2pt, minimum height=8mm, align=center, inner xsep=6pt]
\tikzstyle{big}=[box, minimum width=34mm]
\node[big] (teacher) {Offline semantic teacher\\VLM / human labels / classifier};
\node[big, right=of teacher] (compile) {Predicate compiler\\sparse LUT or int4 linear};
\node[big, right=of compile] (store) {Concept-program store\\$\{F_C\}$};
\draw[->] (teacher) -- node[above]{labels} (compile);
\draw[->] (compile) -- (store);

\node[big, below=14mm of teacher] (item) {Item encoder\\dense embedding $x$};
\node[big, right=of item] (codebox) {Shared low-bit encoding\\1/2/4-bit or binary+LS2};
\node[big, right=of codebox] (code) {Universal item code\\48--192 B/item};
\draw[->] (item) -- (codebox);
\draw[->] (codebox) -- (code);

\node[big, below=14mm of codebox] (query) {Query understanding\\exact + latent predicates};
\node[big, right=of query] (exec) {Semantic execution\\LUT / popcount + AND / NOT};
\node[big, right=of exec] (rank) {Pruned candidates\\downstream ranker};
\draw[->] (query) -- (exec);
\draw[->] (store.south) |- (exec.north);
\draw[->] (code.south) |- (exec.north east);
\draw[->] (exec) -- (rank);
\end{tikzpicture}
\caption{Compiled semantic execution separates expensive supervision from online search. Exact catalog predicates remain conventional filters. Latent predicates are compiled offline into sparse LUT programs or tiny bitwise linear programs over a shared low-bit item representation.}
\label{fig:architecture}
\end{figure*}

Our contributions are fourfold. First, we formulate semantic predicate compilation as a two-sided compression problem: both item bytes and predicate-program bytes matter when a catalog serves many reusable concepts. Second, we study two compiler families over a shared search-side embedding: sparse boosted LUT programs on 1/2/4-bit substrates, and a centered 1-bit two-level reconstruction whose int4 semantic heads execute through bitwise population counts. Third, we retain calibrated log-probability composition as a simple query algebra for independently trained positive and negated predicates. Fourth, we evaluate quality, memory, and native CPU execution against dense retrieval, zero-shot directions, FP32 linear semantics, and PQ64.

The new experiments revise an earlier interpretation of the project. Random 4-bit sparse LUTs remain a strong compressed semantic executor, but they are not the best Pareto point in the current benchmark. The centered binary int4 compiler is smaller on both the item and predicate axes and modestly more accurate, while PQ64 remains the strongest compressed quality and large-pool throughput baseline. We therefore treat random rotation and the particular 4-bit code as design choices, not as the central claim.

We do \emph{not} claim that these codes replace ANN retrieval, nor that the calibrated product rule is a new probabilistic operator. The narrower systems hypothesis is that a search engine can compile expensive semantic supervision into reusable low-bit predicate programs and execute them between broad retrieval and expensive ranking.
